\chapter{Outlook}\label{ch: ch6}%
With the results from the previous sections, it can be concluded that the constructed CV $CV\_dist$ is suitable to drive an MD to replicate the reaction.
This concludes the main goal of this work.
Next, there are many possible directions to further explore the topic of this work.

\subsection{FES}
One goal of a MetaD simulation is to obtain an approximation of the FES. 
Unfortunately, the FES obtained from the MDs driven by the constructed CV was not meaningful and thus, its construction was omitted in this work.
However, as the main goal of the method is to drive the MD from an initial state to a target state, the FES is not crucial for the method.
Nevertheless, it would be interesting to see if a meaningful FES can be obtained with a different construction of the CV.

\subsection{Refined Construction of CVs}
In this work, only one CV ($CV\_dist$) was constructed and used to drive the MD.
However, there are many other metrics that can be used to construct CVs to drive the MD, like the $\| \cdot \|_1$ norm.
Also, the number of PCA components used to construct the CV can be increased to better portray curvilinear motion in the PCA space. 
These adjustments have then to be tested upon the ability to drive an MD and accuracy in replicating the reaction.

\subsection{Alternative MLIPs}
Although MACE-MPA-0 serves as a useful foundation model for materials science, its training data primarily consists of inorganic crystals rather than organic reaction pathways.
Thus, it is not the most accurate MLIP for this usecase. For chemical reactions, other MLIPs e.g. MACE-OFF or ANI-1ccx are more suited \cite{Jacobs2025}.
Therefore, an examination of other MLIPs can be done to see if they can replicate the reaction more accurately and if the constructed CVs can still drive the MD successfully.

\subsection{Application to other systems}
Ultimately, the primary objective of this work is to demonstrate that data driven CVs are able to successfully drive chemical reactions, providing a foundation for automated sampling methods.
The next step would be to apply the method to other systems and reactions to see if it can be generally applied to 
drive MDs from an initial state to a target state for more complex reactions.
Moreover, it would be intriguing to see if the method can also be applied to not only to reactions but also to other systems where a transition from an initial state to a target state is desired.
Some examples of reactions and systems that can be examined are:
\begin{itemize}
    \item other organic reactions
    \item enzymatic reactions
    \item solid state reactions
    \item etc.
\end{itemize}
